\documentclass[12pt]{article}
\usepackage{amsmath}
\usepackage{amssymb}
\usepackage{physics}
\usepackage{graphicx}
\usepackage{hyperref}
\usepackage{booktabs}
\usepackage{array}

\title{The Transformation of Time and Space}
\author{}
\date{}

\begin{document}
\maketitle
 
% ============================================================
\section{Introduction}
% ============================================================

The Lorentz transformation is one of the cornerstones of modern physics. It
relates the measurements of space and time performed by different inertial
observers and underlies virtually every aspect of Special Relativity. Because
of its fundamental importance, it is often introduced early and presented
directly in matrix form. While this approach is efficient, it can leave the
reader with little physical intuition as to why the transformation has the form
it does or what it actually means. \par

In this section we revisit Einstein's original derivation, extending it from
the special case of motion along a coordinate axis to an arbitrary boost
direction. Rather than starting from the final matrix, we derive the Lorentz
transformation from simple physical assumptions, making explicit every step of
the argument. Along the way, we will emphasize an important conceptual point
that is often left implicit: the Lorentz transformation does not determine what
different observers measure, but rather relates the independent measurements
performed in different inertial reference frames.\par

 Einstein original derivation considered the special case in which the relative
 velocity is aligned with the x-axis, so that directions perpendicular to the motion remain unchanged by construction. 
In this section, we follow Einstein's reasoning and extend the derivation to an arbitrary boost direction. Our goals are as follows:

\begin{enumerate}
  \item \textbf{General Boost Transformation:}
    \begin{itemize}
      \item We derive the Lorentz transformation for a general boost, where the relative
            velocity is represented by an arbitrary vector.
      \item This broader approach allows us to handle boosts in any direction, not just along
            the $x$-axis.
    \end{itemize}
  \item \textbf{Transverse Length Invariance:}
    \begin{itemize}
      \item We explicitly discuss the conditions that ensure the preservation of lengths
            perpendicular to the velocity direction.
      \item This consideration is crucial for maintaining consistency across reference frames.
    \end{itemize}
  \item \textbf{Time Transformation via Clock Synchronization:}
    \begin{itemize}
      \item Following Einstein's original method, we will derive the transformation of time
            based on clock synchronization.
      \item This step lays the foundation for our subsequent analysis.
    \end{itemize}
  \item \textbf{Lorentz Transformation for Time and Space:}
    \begin{itemize}
      \item Using the invariance of the speed of light, we obtain compact expressions for the transformations of both space and time.
    \end{itemize}
\end{enumerate}

\textbf{Assumptions and Prerequisites:} We assume that readers are familiar with advanced
calculus, linear algebra, and the fundamental principles of Special Relativity.  \par

\textbf{Notation}. We use Dirac's bra-ket notation as a compact language for vectors, dual vectors, and linear operators. 
Despite its widespread association with quantum mechanics, the notation is purely linear-algebraic 
and is equally applicable to Euclidean and spacetime vector spaces. 
Readers unfamiliar with the notation may consult any standard reference; for a concise introduction, 
see the Wikipedia article on bra-ket notation. \url{https://en.wikipedia.org/wiki/Bra\%E2\%80\%93ket_notation}. \par

Although the Lorentz transformation is often presented directly in matrix form,
deriving it from Einstein's original arguments provides valuable physical
insight into the geometric structure of spacetime and, perhaps more
importantly, into the operational meaning of the transformation itself. Once
this viewpoint is established, many of the central ideas of Special Relativity,
including proper time, four-vectors, and relativistic dynamics, follow
more easily. 

% ============================================================
\section{Synchronization of Clocks}
% ============================================================

In Special Relativity, the concept of clock synchronization plays a fundamental role in
developing the Lorentz transformation of time and space. Although seemingly straightforward,
this concept deserves careful discussion, as defining
synchronization becomes essential for determining when events occur simultaneously.

We begin with a simple thought experiment. Imagine two observers, A and B, both
at rest within the same reference system but positioned at different locations. These
observers, A and B, each have their own clocks, measuring time independently.

Consider an experiment where a photon is emitted from A, travels to B, and gets
reflected back to its origin in A. Observer A records the time of emission as $t_1$, while
observer B notes the time of reflection as $t_2$. Finally, A marks the time of the photon
arrival as $t_3$. We assume that, after the experiment, the two observers compare the times recorded on their respective clocks
(for example they would exchange data after the experiment). A possible outcome is shown in the table below (units of time are not relevant in this
discussion): 

\begin{center}
\begin{tabular}{ccc}
\toprule
Time in A & Time in B & Event \\
\midrule
1 & 2 & Emission, $t_1 = 1$ \\
3 & 4 & Reflection, $t_2 = 4$ \\
5 & 6 & Arrival, $t_3 = 5$ \\
\bottomrule
\end{tabular}
\end{center}

If the observer at point B used the time difference $t_2 - t_1$ to compute the time needed
for the photon to reach B, they would obtain an incorrect result of 3 instead of the actual
value of 2. More worrisome, the time difference $(t_3 - t_2)$ would be 1, falsely suggesting that the two 
observers are not at rest with respect to one another. These
discrepancies occur because the two clocks are not synchronized.

The way for B to synchronize their clock with A is to record the time stamp $t_1$, wait for
the photon to get back to A, ask A for $t_3$, and set the clock in B such that:
\begin{equation}
  t_2 - t_1 = t_3 - t_2
  \label{eq:clock_synchronization}
\end{equation}

which is realized in this case by setting their clock back 1 unit. This requirement defines Einstein synchronization.

Having established a procedure for synchronizing clocks within a single inertial frame, we can now 
investigate how space and time measurements compare between two 
inertial frames in relative motion. This leads directly to the Lorentz transformation. 
This is the subject of the next sections.

% ============================================================
\section{Transformation of Time}
% ============================================================

In this section we will derive the equation for the transformation of time between two
systems of reference moving with constant relative velocity. 

\subsection{The Mirror Experiment}

We now examine the mirror experiment from the perspective of a new reference frame. In this
scenario, the mirror and the emitting source are both at rest within a system of reference
$F'$, which itself is in motion relative to another system of reference $F$. Observers in $F$ assign to each event a time $t$ and a position
$\vb x=(x_1,x_2,x_3)$, while observers in $F'$ assign the corresponding
quantities $t'$ and $\vb x'=(x'_1,x'_2,x'_3)$. We choose the origins of the two reference frames to coincide at 
$t=t'=0$.

\subsubsection*{Experiment Observed from $F'$}

The source is located at the origin of $F'$ and the mirror is positioned at a fixed location in an
 arbitrary direction from the source. A photon is emitted at a time $t'_0$; it reaches the mirror at
time $t'_1$ where it is reflected, and it arrives back at the source at time $t'_2$.

\begin{figure}[h]
  \centering
  \includegraphics[width=0.6\textwidth]{mirror.png}
  \caption{Figure 1: Mirror experiment viewed in the reference frame $F'$. The source and mirror are at rest, while the photon propagates along an arbitrary direction.}
\end{figure}

As we assume that the clocks in $F'$ are synchronized, the following relation holds true:
\begin{equation}
  \frac{1}{2}(t'_0 + t'_2) = t'_1
  \label{eq:clock_synchronization_moving}
\end{equation}

\subsubsection*{Experiment Observed from $F$}

Now we examine the three events---emission, reflection, and arrival back at the source--- from
the viewpoint of reference frame $F$. It is crucial to emphasize that all time, velocity, and
distance quantities in this section \textit{are assumed to be measured within the system of
reference} $F$.

The photon is emitted from the source at time $t_0$ and at the position $\vb{v}\, t_0$ where
$\vb{v}$ is the velocity of $F'$ measured from $F$.

\begin{figure}[h]
  \centering
  \includegraphics[width=0.6\textwidth]{movingframes1.jpg}
  \caption{Figure 2: Velocities and positions of the photon seen from $F$}
\end{figure}

The photon travels to the mirror with speed $\vb{c}_{up} = \vb{c}_1 + \vb{v}$, and it
reaches it at a time:
\[
  t_1 = t_0 + \frac{l}{|\vb{c}_1|} = t_0 + \frac{l}{|\vb{c}_{up} - \vb{v}|}
\]

where $l$ is the distance between the source and the mirror measured
simultaneously in $F$ at the time of emission.
We assume that the source and mirror remain at a fixed distance $l$ from one another throughout the experiment. 
The vector $\vb c_1$ denotes the photon velocity relative to the moving frame $F'$, expressed in the coordinates of $F$.
When the photon reaches the mirror, its position will be
\[
  \vb{v}\, t_1 + (\vb{c}_{up} - \vb{v})(t_1 - t_0)
  = \vb{v}\, t_0 + l\,\frac{\vb{c}_{up}}{|\vb{c}_{up} - \vb{v}|}
\]

The time $t_2$ when the photon reaches the source back again is
\[
  t_2 = t_1 + \frac{l}{|\vb{c}_{down} - \vb{v}|}
      = t_0 + \frac{l}{|\vb{c}_{up} - \vb{v}|} + \frac{l}{|\vb{c}_{down} - \vb{v}|}
\]
where $\vb{c}_{down}$ is the photon velocity returning to the source. \par

The position of the photon when it meets the source is,
\[
  \vb{v}\, t_2
  =
  \vb{v}\, t_0
  +
  \vb{v}\,\frac{l}{|\vb{c}_{up}-\vb{v}|}
  +
  \vb{v}\,\frac{l}{|\vb{c}_{down}-\vb{v}|}.
\]

To derive the Lorentz transformation, we postulate the existence of a
real-valued function on the events described in the reference frame $F$,
\begin{align}
  t' &: F \longrightarrow \mathbb{R}, \\
  t' &= t'(\vb{x},t),
\end{align}
where $(\vb{x},t)$ are the coordinates assigned in $F$ to a given event.

Formally, the events of emission, reflection and meeting the source again will be observed
in $F'$ as
\begin{align}
  t'_0 &= t'(\vb{v}\, t_0,\; t_0) \\
  t'_1 &= t'\!\left(
    \vb{v}\, t_0 + l\,\frac{\vb{c}_{up}}{|\vb{c}_{up} - \vb{v}|},\;
    t_0 + \frac{l}{|\vb{c}_{up} - \vb{v}|}
  \right) \\
  t'_2 &= t'\!\left(
    \vb{v}\, t_0 + \vb{v}\, l
      \bigl(|\vb{c}_{up} - \vb{v}|^{-1} + |\vb{c}_{down} - \vb{v}|^{-1}\bigr),\;
    t_0 + l\bigl(|\vb{c}_{up} - \vb{v}|^{-1} + |\vb{c}_{down} - \vb{v}|^{-1}\bigr)
  \right)
\end{align}
 
Expanding the function
\[
  f(l)
  =
  t'\!\left(
    \vb{v}\,t_0
    +
    l\,\frac{\vb{c}_{up}}
      {|\vb{c}_{up}-\vb{v}|},
    \;
    t_0
    +
    \frac{l}{|\vb{c}_{up}-\vb{v}|}
  \right)
\]
in a Taylor series around $l=0$, we obtain
\[
  t'_1
  =
  t'(\vb{v}\,t_0,t_0)
  +
  l\,|\vb{c}_{up}-\vb{v}|^{-1}
  \ip{
    \grad t'(\vb{v}\,t_0,t_0)
  }{
    (\vb{c}_{up},1)
  }
  +
  o(l).
\]

A similar expansion holds for $t'_2$:
\[
  t'_2
  =
  t'(\vb{v}\,t_0,t_0)
  +
  l\bigl(
    |\vb{c}_{up}-\vb{v}|^{-1}
    +
    |\vb{c}_{down}-\vb{v}|^{-1}
  \bigr)
  \ip{
    \grad t'(\vb{v}\,t_0,t_0)
  }{
    (\vb{v},1)
  }
  +
  o(l).
\]

Using the synchronization condition
\eqref{eq:clock_synchronization_moving} and taking the limit $l\to0$,
we obtain
\[
  \ip{\grad t'}{\vb{w}}=0,
\]
where
\[
  \vb{w}
  =
  \Bigl(
    \vb{v}\bigl(
      |\vb{c}_{up}-\vb{v}|^{-1}
      +
      |\vb{c}_{down}-\vb{v}|^{-1}
    \bigr)
    -
    2\vb{c}_{up}|\vb{c}_{up}-\vb{v}|^{-1},
    \;
    |\vb{c}_{down}-\vb{v}|^{-1}
    -
    |\vb{c}_{up}-\vb{v}|^{-1}
  \Bigr).
\]

Assuming a linear relationship between the time in $F'$ and the coordinates in $F$ as,
\[
  t' = \vb{p} \cdot \vb{x} + \gamma\, t
\]

where $\vb{p}$ is a vector transforming the spatial coordinates and $\gamma$ is the coefficient
transforming time, we arrive at the linear equation in $\vb{p}$ and $\gamma$

\begin{equation}
  \vb{p} \cdot \bigl[
    2\vb{c}_{up}|\vb{c}_{up}-\vb{v}|^{-1}
    -
    \vb{v}\bigl(
      |\vb{c}_{up}-\vb{v}|^{-1}
      +
      |\vb{c}_{down}-\vb{v}|^{-1}
    \bigr)
  \bigr]
  +
  \gamma\bigl(
    |\vb{c}_{up}-\vb{v}|^{-1}
    -
    |\vb{c}_{down}-\vb{v}|^{-1}
  \bigr)
  =0.
  \label{eq:condition_on_p}
\end{equation}

\begin{figure}[h]
  \centering
  \includegraphics[width=0.5\textwidth]{fig1.png}
\caption{Fig. 3. Geometric construction relating the photon velocities and the boost velocity.}
\end{figure}

Equation \eqref{eq:condition_on_p} can be further simplified by examining
the directions of the vectors $\vb v$, $\vb c_{down}$, and
$\vb c_{up}$. Since the photon in $F'$ returns along the same line on
which it traveled toward the mirror, the vectors
$\vb c_{down}-\vb v$ and $\vb c_{up}-\vb v$ are collinear and
oppositely directed, although their magnitudes need not be equal.
Consequently, $\vb c_{up}$, $\vb c_{down}$, and $\vb v$ all lie in the
same plane, as shown in Fig.~3.

In the system of reference of three mutually orthogonal unit vectors $\vb{e}_1$, $\vb{e}_2$,
and $\vb{e}_3$, with $\vb{e}_1$ and $\vb{e}_2$ in the same plane of $\vb{c}_{up}$ and
$\vb{v}$, and $\vb{e}_1$ in the same direction as $\mathbf{v}$, the vector $\mathbf{p}$ has
components, 

\[
  \vb{p} = p_1\,\vb{e}_1 + p_2\,\vb{e}_2 + p_3\,\vb{e}_3
\]

and equation \eqref{eq:condition_on_p} simplifies to

\[
  p_1\, v\bigl(|\vb{c}_{up} - \vb{v}|^{-1} + |\vb{c}_{down} - \vb{v}|^{-1}\bigr)
  - 2c(p_1\cos\alpha + p_2\sin\alpha)|\vb{c}_{up} - \vb{v}|^{-1}
  + \gamma\bigl(|\vb{c}_{up} - \vb{v}|^{-1} - |\vb{c}_{down} - \vb{v}|^{-1}\bigr) = 0
\]

Since the direction of the photon is arbitrary, we choose $\vb e_2$ in
the direction of the component of $\vb p$ orthogonal to $\vb v$. With
this choice, $\vb p$ lies in the plane spanned by $\vb e_1$ and
$\vb e_2$, and therefore $p_3=0$.

From
\[
  \cos\alpha = \frac{c^2 + v^2 - |\vb{c}_{up} - \vb{v}|^2}{2vc}
\]
we arrive at
\[
  p_1\,\frac{c^2}{v} + \gamma
  + p_2\,\frac{2c\sin\alpha\,|\vb{c}_{up} - \vb{v}|}{|\vb{c}_{up} - \vb{v}| - |\vb{c}_{down} - \vb{v}|}
  = 0 \quad \forall\, \alpha, v
\]
which has solutions
\begin{align}
  p_1 &= -\gamma\,\frac{v}{c^2} \\
  p_2 &= 0
\end{align}

The transformation of time between $F'$ and $F$ becomes then
\begin{equation}
  ct' = \gamma\!\left(-\frac{v}{c}\ip{\vb{e}_1}{\vb{x}} + ct\right)
  \label{eq:transformation_of_time}
\end{equation}
where $\gamma$ depends only on the relative velocity between the two frames.

% ============================================================
\section{Transformation of Space}
% ============================================================

\subsection{The Most General Linear Transformation}

The Lorentz transformation describes how the coordinates and the time of the same event
observed in $F$ and $F'$ are related to each other. We postulate that the transformation is linear,

\[
  \vb{X}' = \boldsymbol{\Lambda}(\vb{v})\,\vb{X}
\]

where $\boldsymbol{\Lambda}(\vb{v})$ is a matrix depending only on the velocity vector
$\vb{v}$, and $\vb{X}$ is the four-vector
\[
  \vb{X} = (x_1, x_2, x_3, ct).
\]

A four-vector can be expanded in the basis formed by the mutually
orthogonal unit vectors $\vb{e}_1$, $\vb{e}_2$ and $\vb{e}_3$, with $\vb{e}_1$ chosen parallel to
$\vb{v}$:

\[
  \vb{X} = x_1\ket{\vb{e}_1, 0} + x_2\ket{\vb{e}_2, 0} + x_3\ket{\vb{e}_3, 0}
           + ct\ket{\vb{0}, 1}
\]

and the Lorentz transformation $\boldsymbol{\Lambda}(\vb{v})$ in its most general form
becomes

\begin{equation}
  \begin{aligned}
    \boldsymbol{\Lambda}(\vb{v}) =\;
      &\lambda_{1,1}\ket{\vb{e}_1,0}\bra{\vb{e}_1,0}
      + \lambda_{1,2}\ket{\vb{e}_1,0}\bra{\vb{e}_2,0}
      + \lambda_{1,3}\ket{\vb{e}_1,0}\bra{\vb{e}_3,0}
      + \lambda_{1,4}\ket{\vb{e}_1,0}\bra{\vb{0},1} \\
    +\;&\lambda_{2,1}\ket{\vb{e}_2,0}\bra{\vb{e}_1,0}
      + \lambda_{2,2}\ket{\vb{e}_2,0}\bra{\vb{e}_2,0}
      + \lambda_{2,3}\ket{\vb{e}_2,0}\bra{\vb{e}_3,0}
      + \lambda_{2,4}\ket{\vb{e}_2,0}\bra{\vb{0},1} \\
    +\;&\lambda_{3,1}\ket{\vb{e}_3,0}\bra{\vb{e}_1,0}
      + \lambda_{3,2}\ket{\vb{e}_3,0}\bra{\vb{e}_2,0}
      + \lambda_{3,3}\ket{\vb{e}_3,0}\bra{\vb{e}_3,0}
      + \lambda_{3,4}\ket{\vb{e}_3,0}\bra{\vb{0},1} \\
    +\;&\lambda_{4,1}\ket{\vb{0},1}\bra{\vb{e}_1,0}
      + \lambda_{4,2}\ket{\vb{0},1}\bra{\vb{e}_2,0}
      + \lambda_{4,3}\ket{\vb{0},1}\bra{\vb{e}_3,0}
      + \lambda_{4,4}\ket{\vb{0},1}\bra{\vb{0},1}
  \end{aligned}
  \label{eq:general_Lorentz}
\end{equation}

where $\lambda_{i,j}$ are, at this stage, arbitrary. The transformation of time
\eqref{eq:transformation_of_time} can be utilized to identify some of the $\lambda_{i,j}$ but
not all of them: additional physical assumptions are required on how events are observed in $F$ and $F'$ to
restrict the possible choices for $\lambda_{i,j}$.

\subsection{Determining the Time Transformation}

The transformation $\boldsymbol{\Lambda}(\vb{v})$ applied to a four-vector with only a time
component $\vb{X} = x_4\ket{\vb{0},1}$ gives,
\[
  \boldsymbol{\Lambda}\, x_4\ket{\vb{0},1}
  = x_4\!\left(
      \lambda_{1,4}\ket{\vb{e}_1,0}
    + \lambda_{2,4}\ket{\vb{e}_2,0}
    + \lambda_{3,4}\ket{\vb{e}_3,0}
    + \lambda_{4,4}\ket{\vb{0},1}
  \right)
\]

From the previously derived transformation of time
\eqref{eq:transformation_of_time}, the image of a four-vector having only
a temporal component contains only a component parallel to
$\ket{\vb{e}_1,0}$ and a temporal component. Hence it cannot have
components along $\ket{\vb{e}_2,0}$ or $\ket{\vb{e}_3,0}$, and therefore


\[
\lambda_{2,4}=\lambda_{3,4}=0.
\]

\subsection{First Assumption: Preservation of the Longitudinal Direction}

The first assumption we make is that the subspace spanned by the direction
$\vb e_1$ and the time direction is invariant. Thus, an event in $F$ with
spatial components only along $\vb e_1$ will not acquire components along
$\vb e_2$ or $\vb e_3$ when observed in $F'$. Given an event

$\vb{X} = x_1\ket{\vb{e}_1, 0} + x_4\ket{\vb{0}, 1}$, its transformation in $F'$ will be,
\[
  \begin{aligned}
    \boldsymbol{\Lambda}(\vb{v})\bigl(x_1\ket{\vb{e}_1,0} + x_4\ket{\vb{0},1}\bigr)
    &= x_1\!\left(
        \lambda_{1,1}\ket{\vb{e}_1,0}
      + \lambda_{2,1}\ket{\vb{e}_2,0}
      + \lambda_{3,1}\ket{\vb{e}_3,0}
      + \lambda_{4,1}\ket{\vb{0},1}
    \right) \\
    &\quad + x_4\!\left(\lambda_{1,4}\ket{\vb{e}_1,0} + \lambda_{4,4}\ket{\vb{0},1}\right)
  \end{aligned}
\]

Since this condition must hold for every value of $x_1$, the coefficients
of the basis vectors $\ket{\vb{e}_2,0}$ and $\ket{\vb{e}_3,0}$ must vanish.
Hence,
\[
  \lambda_{2,1}=\lambda_{3,1}=0.
\]

The general form \eqref{eq:general_Lorentz} further simplifies as,
\[
  \begin{aligned}
    \boldsymbol{\Lambda}(\vb{v}) =\;
      &\lambda_{1,1}\ket{\vb{e}_1,0}\bra{\vb{e}_1,0}
      + \lambda_{1,2}\ket{\vb{e}_1,0}\bra{\vb{e}_2,0}
      + \lambda_{1,3}\ket{\vb{e}_1,0}\bra{\vb{e}_3,0}
      + \lambda_{1,4}\ket{\vb{e}_1,0}\bra{\vb{0},1} \\
    +\;&\lambda_{2,2}\ket{\vb{e}_2,0}\bra{\vb{e}_2,0}
      + \lambda_{2,3}\ket{\vb{e}_2,0}\bra{\vb{e}_3,0} \\
    +\;&\lambda_{3,2}\ket{\vb{e}_3,0}\bra{\vb{e}_2,0}
      + \lambda_{3,3}\ket{\vb{e}_3,0}\bra{\vb{e}_3,0} \\
    +\;&\lambda_{4,1}\ket{\vb{0},1}\bra{\vb{e}_1,0}
      + \lambda_{4,2}\ket{\vb{0},1}\bra{\vb{e}_2,0}
      + \lambda_{4,3}\ket{\vb{0},1}\bra{\vb{e}_3,0}
      + \lambda_{4,4}\ket{\vb{0},1}\bra{\vb{0},1}
  \end{aligned}
\]

\subsection{Second Assumption: Invariance Orthogonal to the Boost Direction}

The second assumption needed is that the spatial components along $\ket{\vb{e}_2,0}$ and
$\ket{\vb{e}_3,0}$ of any event preserve their direction and length under
$\boldsymbol{\Lambda}(\vb{v})$. 

For an event $\vb{X} = x_2\ket{\vb{e}_2, 0} + x_3\ket{\vb{e}_3, 0}$, the corresponding event in $F'$ will be,

\[
  \begin{aligned}
    \boldsymbol{\Lambda}(\vb{v})
    \bigl(
      x_2\ket{\vb{e}_2,0}
      +
      x_3\ket{\vb{e}_3,0}
    \bigr)
    &=
    x_2\!\left(
        \lambda_{1,2}\ket{\vb{e}_1,0}
      + \lambda_{2,2}\ket{\vb{e}_2,0}
      + \lambda_{3,2}\ket{\vb{e}_3,0}
      + \lambda_{4,2}\ket{\vb{0},1}
    \right)
    \\
    &\quad
    +
    x_3\!\left(
        \lambda_{1,3}\ket{\vb{e}_1,0}
      + \lambda_{2,3}\ket{\vb{e}_2,0}
      + \lambda_{3,3}\ket{\vb{e}_3,0}
      + \lambda_{4,3}\ket{\vb{0},1}
    \right).
  \end{aligned}
\]

Our assumption therefore requires $\lambda_{1,2} = \lambda_{1,3} = \lambda_{2,3} = \lambda_{3,2}
= 0$ and $\lambda_{2,2} = \lambda_{3,3} = 1$, which results in a further simplification of
$\boldsymbol{\Lambda}(\vb{v})$:
\[
  \begin{aligned}
    \boldsymbol{\Lambda}(\vb{v}) =\;
      &\lambda_{1,1}\ket{\vb{e}_1,0}\bra{\vb{e}_1,0}
      + \lambda_{1,4}\ket{\vb{e}_1,0}\bra{\vb{0},1} \\
    +\;&\ket{\vb{e}_2,0}\bra{\vb{e}_2,0} \\
    +\;&\ket{\vb{e}_3,0}\bra{\vb{e}_3,0} \\
    +\;&\lambda_{4,1}\ket{\vb{0},1}\bra{\vb{e}_1,0}
      + \lambda_{4,2}\ket{\vb{0},1}\bra{\vb{e}_2,0}
      + \lambda_{4,3}\ket{\vb{0},1}\bra{\vb{e}_3,0}
      + \lambda_{4,4}\ket{\vb{0},1}\bra{\vb{0},1}
  \end{aligned}
\]

Using the previously derived  \eqref{eq:transformation_of_time}, we identify

\[
  \lambda_{4,1}=-\gamma\frac{v}{c},
  \qquad
  \lambda_{4,2}=\lambda_{4,3}=0,
  \qquad
  \lambda_{4,4}=\gamma.
\]

We therefore arrive at the final functional form of
$\boldsymbol{\Lambda}(\vb{v})$:

\begin{equation}
  \begin{aligned}
    \boldsymbol{\Lambda}(\vb{v}) =\;
      &\lambda_{1,1}\ket{\vb{e}_1,0}\bra{\vb{e}_1,0}
      + \lambda_{1,4}\ket{\vb{e}_1,0}\bra{\vb{0},1}
      + \ket{\vb{e}_2,0}\bra{\vb{e}_2,0}
      + \ket{\vb{e}_3,0}\bra{\vb{e}_3,0} \\
      &- \gamma\,\frac{v}{c}\ket{\vb{0},1}\bra{\vb{e}_1,0}
      + \gamma\ket{\vb{0},1}\bra{\vb{0},1}
  \end{aligned}
  \label{eq:simplified_Lorentz}
\end{equation}

Finally, the coefficients $\lambda_{1,1}$ and $\gamma$ must be positive
to preserve the orientation of the longitudinal spatial axis and the direction of increasing time.

\subsection{Third Assumption: Invariance of Light Speed}

In Special Relativity the spacetime interval between two events $(x_1, y_1, z_1, ct_1)$ and
$(x_2, y_2, z_2, ct_2)$ is defined as

\[
  \Delta s^2 = (x_2 - x_1)^2 + (y_2 - y_1)^2 + (z_2 - z_1)^2 - c^2(t_2 - t_1)^2
\]

and can be written in bi-linear form as

\[
  \Delta s^2 = \ip{\vb{X}}{\vb{M}\,\vb{X}}
\]

where $\vb{X} = (x_2 - x_1,\; y_2 - y_1,\; z_2 - z_1,\; c(t_2 - t_1))$ and
\[
  \vb{M} =
  \begin{pmatrix}
    1 & 0 & 0 &  0 \\
    0 & 1 & 0 &  0 \\
    0 & 0 & 1 &  0 \\
    0 & 0 & 0 & -1
  \end{pmatrix}
\]

A spacetime interval measured in $F'$ is linked to the same interval in $F$ by the Lorentz
transformation

\[
  \Delta s'^2
  = \ip{\vb{X}'}{\vb{M}\,\vb{X}'}
  = \ip{\boldsymbol{\Lambda}\vb{X}}{\vb{M}\,\boldsymbol{\Lambda}\vb{X}}
  = \ip{\vb{X}}{\boldsymbol{\Lambda}^T\vb{M}\,\boldsymbol{\Lambda}\,\vb{X}}
\]

In the system $F$ we consider a photon moving along a unit vector $\vb{n}$. The photon at
time $t_1$ is in the position $\vb{x}_1$, and moves to the position
$\vb{x}_1 + \vb{n}\,c(t_2 - t_1)$ at the time $t_2$. 

Since $||\vb n||=1$ the spacetime between the two events is identically zero,

\[
  \Delta s^2 = c^2(t_2 - t_1)^2\|\vb{n}\|^2 - c^2(t_2 - t_1)^2 = 0
\]


When we examine two events occurring along the trajectory of a photon (referred to as the
photon's world line), the spacetime interval $\Delta' s^2$ in $F'$ must also be zero, assuming
that the speed of light $c$ remains invariant,
\[
  \ip{
    \bigl(c(t_2-t_1)\vb{n},\; c(t_2-t_1)\bigr)
  }{
    \boldsymbol{\Lambda}^T\vb{M}\,\boldsymbol{\Lambda}
    \cdot\bigl(c(t_2-t_1)\vb{n},\; c(t_2-t_1)\bigr)
  } = 0
\]
equivalent to,
\begin{equation}
  \ip{(\vb{n},1)}{\boldsymbol{\Lambda}^T\vb{M}\,\boldsymbol{\Lambda}\cdot(\vb{n},1)}
  = 0
  \label{eq:invariance_light_speed}
\end{equation}
for all $\vb{n}$ such that $\|\vb{n}\| = 1$.

To solve the equation above, we need $\vb{M}\,\boldsymbol{\Lambda}$, which is
$\boldsymbol{\Lambda}$ with the time component's sign changed,
\[
  \begin{aligned}
    \vb{M}\,\boldsymbol{\Lambda}(\vb{v}) =\;
      &\lambda_{1,1}\ket{\vb{e}_1,0}\bra{\vb{e}_1,0}
      + \lambda_{1,4}\ket{\vb{e}_1,0}\bra{\vb{0},1}
      + \lambda_{2,2}\ket{\vb{e}_2,0}\bra{\vb{e}_2,0}
      + \lambda_{3,3}\ket{\vb{e}_3,0}\bra{\vb{e}_3,0} \\
      &+ \gamma\,\frac{v}{c}\ket{\vb{0},1}\bra{\vb{e}_1,0}
      - \gamma\ket{\vb{0},1}\bra{\vb{0},1}
  \end{aligned}
\]

and $\boldsymbol{\Lambda}^T$ is readily obtained by interchanging the bra and ket vectors,
\[
  \begin{aligned}
    \boldsymbol{\Lambda}(\vb{v})^T =\;
      &\lambda_{1,1}\ket{\vb{e}_1,0}\bra{\vb{e}_1,0}
      + \lambda_{1,4}\ket{\vb{0},1}\bra{\vb{e}_1,0}
      + \lambda_{2,2}\ket{\vb{e}_2,0}\bra{\vb{e}_2,0}
      + \lambda_{3,3}\ket{\vb{e}_3,0}\bra{\vb{e}_3,0} \\
      &- \gamma\,\frac{v}{c}\ket{\vb{e}_1,0}\bra{\vb{0},1}
      + \gamma\ket{\vb{0},1}\bra{\vb{0},1}
  \end{aligned}
\]



The operator $\boldsymbol{\Lambda}(\vb{v})^T\vb{M}\,\boldsymbol{\Lambda}(\vb{v})$ is
\[
  \begin{aligned}
    &\left(\lambda_{1,1}^2 - \gamma^2\frac{v^2}{c^2}\right)\ket{\vb{e}_1,0}\bra{\vb{e}_1,0} \\
    +\;&\left(\lambda_{1,1}\lambda_{1,4} + \gamma^2\frac{v}{c}\right)
      \bigl(\ket{\vb{0},1}\bra{\vb{e}_1,0} + \ket{\vb{e}_1,0}\bra{\vb{0},1}\bigr) \\
    +\;&\lambda_{2,2}^2\ket{\vb{e}_2,0}\bra{\vb{e}_2,0}
      + \lambda_{3,3}^2\ket{\vb{e}_3,0}\bra{\vb{e}_3,0} \\
    +\;&\left(\lambda_{1,4}^2 - \gamma^2\right)\ket{\vb{0},1}\bra{\vb{0},1}
  \end{aligned}
\]

and the equation for the invariance of speed of light \eqref{eq:invariance_light_speed} reads,
\[
  \left(\lambda_{1,1}^2 - \gamma^2\frac{v^2}{c^2}\right)n_1^2
  + 2\left(\lambda_{1,1}\lambda_{1,4} + \gamma^2\frac{v}{c}\right)n_1
  + n_2^2 + n_3^2
  + \left(\lambda_{1,4}^2 - \gamma^2\right) = 0
\]

The above equation is identically zero for all $\vb{n}$ such that $\|\vb{n}\| = 1$ only if,

\begin{align}
  \lambda_{1,1}^2 - \gamma^2\displaystyle\frac{v^2}{c^2} &= 1 \\
  \lambda_{1,1}\lambda_{1,4} + \gamma^2\displaystyle\frac{v}{c} &= 0 \\
  \lambda_{1,4}^2 - \gamma^2 &= -1
\end{align}

which admits the following solution, continuously connected to the
identity transformation at $v=0$:
\begin{align}
  \gamma
  &=
  \frac{1}{\sqrt{1-\frac{v^2}{c^2}}},
  \\
  \lambda_{1,1}
  &=
  \gamma,
  \\
  \lambda_{1,4}
  &=
  -\gamma\frac{v}{c}.
\end{align}

The Lorentz transformation takes the final form

\begin{equation}
  \begin{aligned}
    \boldsymbol{\Lambda}(\vb{v}) =\;
      &\gamma\ket{\vb{e}_1,0}\bra{\vb{e}_1,0}
      - \gamma\,\frac{v}{c}
        \ket{\vb{e}_1,0}\bra{\vb{0},1}
      + \ket{\vb{e}_2,0}\bra{\vb{e}_2,0}
      + \ket{\vb{e}_3,0}\bra{\vb{e}_3,0}
      \\
      &- \gamma\,\frac{v}{c}
        \ket{\vb{0},1}\bra{\vb{e}_1,0}
      + \gamma\ket{\vb{0},1}\bra{\vb{0},1}.
  \end{aligned}
  \label{eq:Lorentz_transformation}
\end{equation}

The inverse of $\boldsymbol{\Lambda}(\vb{v})$ transforms coordinates from
$F'$ to $F$. It is obtained by reversing the boost velocity and can be
verified directly:

\begin{equation}
  \boldsymbol{\Lambda}^{-1}(\vb{v}) = \boldsymbol{\Lambda}(-\vb{v}) =
    \gamma\ket{\vb{e}_1,0}\bra{\vb{e}_1,0}
    + \gamma\,\frac{v}{c}\ket{\vb{e}_1,0}\bra{\vb{0},1}
    + \ket{\vb{e}_2,0}\bra{\vb{e}_2,0}
    + \ket{\vb{e}_3,0}\bra{\vb{e}_3,0}
    + \gamma\,\frac{v}{c}\ket{\vb{0},1}\bra{\vb{e}_1,0}
    + \gamma\ket{\vb{0},1}\bra{\vb{0},1}
  \label{eq:Lorentz_inverse}
\end{equation}

The Lorentz transformation written explicitly in terms of the velocity vector, using
$\ket{\vb{e}_1,0} = \ket{\vb{v},0}/v$, reads,
\begin{equation}
  \boldsymbol{\Lambda}(\vb{v}) =
    \frac{\gamma}{v^2}\ket{\vb{v},0}\bra{\vb{v},0}
    - \frac{\gamma}{c}\ket{\vb{v},0}\bra{\vb{0},1}
    + \ket{\vb{e}_2,0}\bra{\vb{e}_2,0}
    + \ket{\vb{e}_3,0}\bra{\vb{e}_3,0}
    - \frac{\gamma}{c}\ket{\vb{0},1}\bra{\vb{v},0}
    + \gamma\ket{\vb{0},1}\bra{\vb{0},1}
  \label{eq:Lorentz_in_velocity}
\end{equation}

When $\vb{v}$ has the direction of the $x$-axis, the transformation simplifies to the original
Einstein formula:
\begin{align}
  x' &= \gamma(x - vt)               \nonumber \\
  y' &= y                            \nonumber \\
  z' &= z                            \nonumber \\
  t' &= \gamma\!\left(t - \frac{v}{c^2}x\right)
  \label{eq:Lorentz_1d}
\end{align}

Finally, one readily verifies that the Lorentz transformation satisfies the condition
$\boldsymbol{\Lambda}^T\vb{M}\,\boldsymbol{\Lambda} = \vb{M}$, hence $\Delta s'^2=\Delta s^2$ for every pair of events, not only for
events lying on the worldline of a photon:

\begin{equation}
  \Delta s'^2
  = \ip{\vb{X}'}{\vb{M}\,\vb{X}'}
  = \ip{\boldsymbol{\Lambda}\vb{X}}
        {\vb{M}\,\boldsymbol{\Lambda}\vb{X}}
  = \ip{\vb{X}}
        {\boldsymbol{\Lambda}^T\vb{M}\,
         \boldsymbol{\Lambda}\,\vb{X}}
  = \ip{\vb{X}}{\vb{M}\,\vb{X}}
  = \Delta s^2.
  \label{eq:spacetime_interval_invariance}
\end{equation}

\section{A Final Remark on the Meaning of the Lorentz Transformation}

When studying the Lorentz transformation, it is important to keep in mind a simple—perhaps 
obvious—but surprisingly often unstated point.\par

Each inertial reference frame possesses its own system of clocks and rulers.
Observers at rest in the frame $F$ measure space and time using their own clocks
and rulers, while observers at rest in $F'$ perform the same measurements using
their own independent instruments. The uniform motion of one frame relative to
the other does not physically modify either system. This principle was already
implicit in Galilean relativity and remains equally valid in Special Relativity.\par

Imagine that we remain on Earth, going about our everyday lives, while an alien
spacecraft flies past us at a speed approaching that of light. Suppose that the
alien civilization is familiar with our way of measuring time. Through previous
encounters with humanity, they know what we mean by a second, a minute, an
hour, and a day.\par

Nothing unusual would happen from our perspective when the aliens fly by. Our
clocks would continue to tick at their normal rate, our rulers would retain
their usual lengths, and the laws of physics would proceed exactly as they
always have. To us, there is nothing remarkable about our relative motion with
respect to the alien spacecraft. In fact, we would not even know that an alien
spacecraft was observing our motion. We would wake up, go to work,
build telescopes, and perform experiments exactly as we always have. Our clocks
would continue to tick normally, our rulers would retain their usual lengths,
and the space and time we measure would simply describe the universe in which
we live.\par

The situation is exactly the same for the crew aboard the alien spacecraft.
They, too, carry their own clocks and rulers, and from their perspective it is
our planet that is rushing past at enormous speed. Their clocks tick normally,
their rulers do not contract, and the physical world they experience is every
bit as ordinary as ours. Each observer naturally believes that their own clocks
and rulers describe the world correctly, and they are both right.\par

The alien observers, however, would describe our world very differently. They
would conclude that one of our days lasts much longer than twenty-four hours,
and that our cars, airplanes, and every object moving with the Earth are
contracted along the direction of motion, as we shall derive in the following
chapters.\par

What the alien observers \textbf{can} do is make sense of these apparently very
long days by using the Lorentz transformation. Since they know the relative
velocity between our two reference frames, they can compute how long our days
actually last \textbf{for us}. They would then conclude that one day measured by
our own clocks --- our \emph{proper time} --- is exactly twenty-four hours, even though they \textbf{measure} a much longer time interval
\emph{for} the same process in their own reference frame. The Lorentz transformation therefore allows them to translate
their description of our measurements into the measurements made by our own
clocks and rulers. \par

It is important to appreciate the role of the Lorentz transformation in this
process. The alien observers never directly measure the passage of time on our
clocks. What they measure is the duration of one of our days according to
\emph{their} own clocks. Likewise, we never directly measure the passage of
time on theirs. Every observer performs measurements using only the clocks and
rulers belonging to their own reference frame.

The purpose of the Lorentz transformation is therefore not to determine what an
observer measures. Those measurements are obtained directly by experiment.
Rather, its role is to reconcile the independent measurements performed in the
two reference frames, allowing each observer to infer what the other observer
measures using their own clocks and rulers. \par

Interestingly, if one of us were able to observe the alien spacecraft as it
passed by, we would describe their world in exactly the same way. We would
conclude that their days are much longer than ours and that distances along the
direction of motion are contracted. Each observer therefore regards the other
reference frame as exhibiting time dilation and length contraction, yet both
descriptions are equally valid. \par

This is a profound departure from our own experience. Contrary to the
Newtonian view, there exists no preferred or absolute frame of reference against
which motion can be measured. There exists no hidden or absolute reference frame from which one could 
decide whether our description or that of the alien observers is the "correct" one. 
\par

The Lorentz transformation should therefore not be interpreted as imposing a
relationship between the physical clocks or rulers belonging to the two
reference frames. Rather, it provides the mathematical rule that relates the
coordinates assigned by the two observers to the \emph{same spacetime event}. 
The transformation does not constrain the
independent evolution of time or space within either frame; it simply
relates two different descriptions of one physical reality. \par

This distinction is essential. The Lorentz transformation is fundamentally a
transformation between coordinate systems, not a law describing how clocks
or rulers physically change because of motion. The phenomena of time
dilation, length contraction, and the relativity of simultaneity arise only
when comparing measurements of specific physical processes performed in
different inertial frames through these coordinate transformations. \par

 
There is, however, one final viewpoint that deserves special attention. Besides the reference frames
$F$ and $F'$, every physical object carries with it its own rest frame. This
frame is not privileged in any universal sense, but it is privileged for the
object itself: it is the frame in which the object remains at rest and in which
its own clocks measure the passage of its history. The time measured by these
clocks is called the \emph{proper time}. Proper time is one of the fundamental quantities of Special Relativity and 
will become the natural parameter describing the history of every physical system throughout the chapters that follow. \par

It is worth emphasizing one final conceptual point. \par

Every observer performs measurements using only the clocks and rulers at rest
in their own reference frame. The Lorentz transformation does not determine
those measurements; they are obtained directly by experiment. Its purpose is to
relate the independent measurements performed in different inertial reference
frames. \par

\end{document}
